\documentclass{report}
\usepackage{amsmath,amssymb}
\setlength{\parindent}{0mm}
\setlength{\parskip}{1em}
\begin{document}
\begin{center}
\rule{6in}{1pt} \
{\large Yury Gryazin \\
{\bf Simplified Krylov subspace algorithm for a high-order compact approximation of 3-D Helmholtz equation.}}

Mathematics Department \\ Idaho State University \\ Stop 8085 \\ Pocatello ID 83209
\\
{\tt gryazin@isu.edu}\end{center}

In recent years, the problem of increasing the resolution of existing
numerical solvers has become an urgent task in many areas of science and
engineering. Most of the existing efficient solvers for structured
matrices were developed for lower-order approximation discretized partial
differential equations. The need for improved accuracy of underlying
algorithms leads to modified discretized systems and as a result to the
modification of the numerical solvers. As an example, one can consider a
modification of Fast Fourier Transform (FFT) type algorithms for the
solution of the Helmholtz equation when the accuracy of the
discretization of the original equation changes from a second to a fourth
order approximation. In many situations, this modification has been
successfully implemented but it usually has taken a major revision of the
underlying numerical algorithm.
In this paper, we consider the use of the lower-order approximation
solver as a preconditioner for an efficient implementation of a
high-order approximation scheme for the Helmholtz equation in the Krylov
subspace method framework. As a first step, we develop a compact
sixth-order approximation finite-difference scheme for the 3D Helmholtz
equation with Dirichlet, Neumann and Sommerfeld-like boundary conditions.
Then we consider the solution of the proposed discretized system by using
the Generalized Minimal Residual (GMRES) method with a second order
approximation system as a preconditioner. The numerical implementation of
this approach requires the solution of the low-order approximation system
on each step of the iterative algorithm, which can be efficiently done by
a FFT-type method.
The convergence analysis of the GMRES method in the case of a 1-D
preconditioned system suggests a simplified Krylov subspace method which
theoretically converges slower than the original GMRES algorithm but that
does not require the Arnoldi orthogonalization procedure. We will use the
abbreviation SKSM to describe this algorithm.
In our numerical experiments, we demonstrate the accuracy of the proposed
compact six-order approximation scheme and the efficiency of the proposed
iterative procedures based on the combination of GMRES-FFT and SKSM-FFT.
Both methods exhibit excellent convergence in most considered test
problems and can be viewed as general blueprints for the improvement of
the accuracy of existing lower-order approximation solvers. We also
present the results of a numerical experiment which demonstrates that in
some situations the GMRES-FFT combination experiences stagnation and
needs to be restarted, while the SKSM-FFT method continues to exhibit
excellent convergence properties.


\end{document}
